\documentclass{reflection}

\fieldwork{1}
\topic{Uniformly accelerated motion}
\duedate{Monday, September 8}

\begin{document}
\maketitle

Find an example of the following:
\begin{quote}
    A roughly constant force other than gravity which, without changing the force, brings an object first to a stop and then accelerates it in the opposite direction. Be sure to address the following:
    \begin{enumerate}
        \item What is the source of the force on the object?
        \item What is the `shape' of the motion?
        \item What is one thing about this motion that is not consistent with your understanding of uniformly accelerated motion in one dimension? What might be the cause of this discrepancy?
    \end{enumerate}
\end{quote}

To receive full credit, this must be an example which you reasonably might have experienced in the course of this week. Also, it must be correct as an example of the above phenomenon.

\ruledlines{19}

\end{document}
